
\section{关于 \texorpdfstring{$dF$}{dF} 的两个引理}
\label{sec:two-lemmas-df}

我们要证明下列两条引理。

\emph{引理（引理~9.4.8）。}设 $r_{0}>0$。存在一个常数 $K>0$，使得对每个
$\rho\ge\rho_{0}$ 与每个 $Z\in\overline{B}(0,r_{0})\subset W^{1,p}(\RR\times S^{1};\RR^{2n})$，我们都有
\[
\bigl\|(dF_{\rho})_{Z}-(dF_{\rho})_{0}\bigr\|^{\mathrm{op}}\le K\|Z\|_{W^{1,p}}\text{。}
\]

\emph{引理（引理~9.4.16）。}设 $r_{0}>0$。存在一个常数 $M>0$（它不依赖于 $\rho$），使得对每个 $\rho\ge\rho_{0}$ 与每个 $Z\in\overline{B}(0,r_{0})\subset W^{1,p}(\RR\times S^{1};\RR^{2n})$，我们都有
\[
\Bigl\|\frac{\partial F}{\partial\rho}(\rho,Z)-\frac{\partial F}{\partial\rho}(\rho,0)\Bigr\|_{L^{p}}\le M\|Z\|_{W^{1,p}}\text{。}
\]

\begin{remark}\label{thm:op-norm-remark}
在引理~9.4.8 的陈述中，$\|\cdot\|^{\mathrm{op}}$ 表示算子范数。一条蕴含这两条引理的断言是：$F$ 是 $\cC^{2}$ 类的，而且算子范数 $\|d^{2}F\|^{\mathrm{op}}$ 在 $[\rho_{0},+\infty[\,\times\,\overline{B}(0,r_{0})$ 上\trBounded。
\end{remark}

我们从证明中将要用到的若干初等结果开始。首先有\cref{thm:w1p-product-inequality}（它断言：一个 $W^{1,p}$ 函数与一个\trBounded{} $\cC^{\infty}$ 函数的乘积是一个 $W^{1,p}$ 函数）。下面两条引理中的第一条就属于同一类型。

\begin{lemma}\label{thm:w1p-product-lp}
设 $p>2$，并设 $f$ 与 $g$ 是 $W^{1,p}(\RR\times S^{1};\RR)$ 的两个元素。那么存在一个常数 $K>0$，使得
\[
\|fg\|_{L^{p}}\le K\|f\|_{L^{p}}\|g\|_{W^{1,p}}\text{。}
\]
\end{lemma}

\begin{proof}
$p$ 大于 $2$ 这一事实使我们得以使用连续\trInjective
\[
W^{1,p}(\RR\times S^{1};\RR)\longrightarrow L^{\infty}(\RR\times S^{1};\RR)
\]
它具有性质：对某个常数 $K>0$，$\|h\|_{\infty}\le K\|h\|_{W^{1,p}}$。于是我们得到
\[
\begin{aligned}
\left(\int_{\RR\times S^{1}}|fg|^{p}\right)^{1/p}
&\le\left(\int_{\RR\times S^{1}}|f|^{p}\,\|g\|_{\infty}^{p}\right)^{1/p}\\
&\le\|g\|_{\infty}\|f\|_{L^{p}}\le K\|f\|_{L^{p}}\|g\|_{W^{1,p}}\text{。}
\end{aligned}
\]
\end{proof}

\begin{lemma}\label{thm:a-b-estimates}
设 $W$ 是一个\trCompact\trManifold，并设
\[
A\colon W\times\RR^{m}\times\RR^{m}\longrightarrow\RR^{m}\text{，}\qquad
B\colon W\times\RR^{m}\times\RR^{m}\times\RR^{m}\longrightarrow\RR^{m}
\]
是两个 $\cC^{\infty}$ 映射。我们假设：对每个 $p\in W$ 与每个 $X\in\RR^{m}$，$A(p,X,\cdot)$ 是线性的，$B(p,X,\cdot,\cdot)$ 是双线性的。那么对每个 $r_{1}>0$，都存在一个常数 $C>0$，使得对每个 $p\in W$ 与每个 $X\in\overline{B}(0,r_{1})\subset\RR^{m}$，我们有：
\begin{enumerate}[label=(\arabic*)]
\item 对每个 $v\in\RR^{m}$，$\|A(p,X,v)\|\le C\|v\|$。
\item 对所有 $v_{1},v_{2}\in\RR^{m}$，
\[
\|A(p,X,v_{2})-A(p,0,v_{1})\|\le C\left(\|X\|\|v_{1}\|+\|v_{2}-v_{1}\|\right)\text{。}
\]
\item 对所有 $v,Y\in\RR^{m}$，$\|B(p,X,v,Y)\|\le C\|v\|\|Y\|$。
\item 对所有 $v,Y_{1},Y_{2}\in\RR^{m}$，
\[
\|B(p,X,v,Y_{2})-B(p,0,v,Y_{1})\|\le C\left(\|X\|\|v\|\|Y_{1}\|+\|v\|\|Y_{2}-Y_{1}\|\right)\text{。}
\]
\end{enumerate}
\end{lemma}

\begin{proof}
这个证明是直接而初等的：
\begin{enumerate}[label=(\arabic*)]
\item 我们把 $A$ 视为一个映射
\[
A\colon W\times\RR^{m}\longrightarrow\scr{L}(\RR^{m},\RR^{m})\text{。}
\]
由于它是连续的，\trCompact{}集 $W\times\overline{B}(0,r_{1})$ 的\trImage{}是\trCompact{}的，从而按算子范数\trBounded{}，这就证明了第一条不等式。
\item 由于 $A$ 是 $\cC^{1}$ 类的，我们有
\[
\|A(p,X)-A(p,0)\|^{\mathrm{op}}\le\sup_{t\in[0,1]}\bigl\|(dA)_{p,tX}\bigr\|\,\|X\|\le C\|X\|\text{。}
\]
于是我们得到
\[
\begin{aligned}
\|A(p,X,v_{2})-A(p,0,v_{1})\|&\le\|A(p,X,v_{2}-v_{1})\|+\|A(p,X,v_{1})-A(p,0,v_{1})\|\\
&\le C\|v_{2}-v_{1}\|+\|A(p,X)-A(p,0)\|^{\mathrm{op}}\|v_{1}\|\\
&\text{（用 (1)）}\\
&\le C\|v_{2}-v_{1}\|+C\|X\|\|v_{1}\|\text{。}
\end{aligned}
\]
\item 我们把 $B$ 视为一个从 $W\times\RR^{m}$ 到 $\RR^{m}$ 上的双线性映射的映射，并照 (1) 那样进行。
\item 我们写出
\[
\begin{aligned}
&\|B(p,X,v,Y_{2})-B(p,0,v,Y_{1})\|\\
&\qquad\le\|B(p,X,v,Y_{2}-Y_{1})\|+\|B(p,X,v,Y_{1})-B(p,0,v,Y_{1})\|
\end{aligned}
\]
并照 (2) 那样进行。
\end{enumerate}
\end{proof}
